\title{The Era of 1-bit LLMs: All Large Language Models are in 1.58 Bits}

\begin{abstract}
Emerging advancements like BitNet~\cite{bitnet} are ushering in a transformative period for 1-bit Large Language Models (LLMs). Here, we propose a novel 1-bit LLM architecture, termed \textbf{\text{BitNet b1.58}}, wherein each parameter or weight of the model is drawn from a ternary set \{-1, 0, 1\}. This design attains perplexity and downstream task outcomes comparable to conventional full-precision (FP16 or BF16) Transformer LLMs of identical model scale and exposure to training data, yet delivers substantial improvements in latency, memory footprint, throughput, and energy efficiency. More notably, the 1.58-bit LLM establishes an innovative scaling principle and provides guidelines for building subsequent generations of LLMs that do not compromise between performance and efficiency. Additionally, this work enables the development of new computational paradigms and encourages the creation of custom hardware specifically tailored for 1-bit LLMs.
\end{abstract}

\section{The Era of 1-bit LLMs}

Artificial Intelligence research has recently experienced significant advancements in both the size and capability of Large Language Models (LLMs). While these models have achieved outstanding results across numerous natural language processing benchmarks, their escalating complexity has resulted in deployment difficulties, coupled with rising environmental and financial costs owing to elevated energy demands.

To mitigate these issues, post-training quantization has been explored as a method to compress model precision for inference~\cite{smoothquant, gptq, quip, quip_sharp}. By reducing the bit-width of both weights and activations, this approach yields considerable savings in terms of memory usage and computation for LLMs. Industry trends have moved from 16-bit formats down to compact representations, including 4-bit models~\cite{gptq, awq}. Nevertheless, despite widespread industrial adoption, post-training quantization remains a sub-optimal solution.

Advancements in 1-bit model architectures, exemplified by BitNet~\cite{bitnet}, indicate a viable path for decreasing the operational cost of LLMs without sacrificing their efficacy. Standard LLMs utilize 16-bit floating point numbers (FP16 or BF16), and their computational load is predominantly due to extensive matrix multiplications. The primary expense arises from performing floating-point multiplications and additions. In comparison, BitNet’s matrix multiplication operations depend solely on integer addition, thereby delivering a substantial reduction in energy consumption for LLMs. Since power serves as a key limiter to computational throughput in most hardware accelerators, the energy gains offered by such approaches can also result in faster computation.

Beyond compute, a significant expense during inference is associated with moving model weights from DRAM to an accelerator’s local memory (e.g., SRAM). While increasing SRAM can improve bandwidth and throughput, the associated cost far exceeds that of DRAM. 1-bit LLMs offer a drastically smaller memory requirement, both in terms of storage capacity and transfer bandwidth, as compared to their full-precision counterparts. This attribute can notably minimize both the expense and latency of transferring model weights from DRAM, resulting in more rapid and cost-effective inference.

Here, we propose a notable 1-bit LLM extension, \textbf{\text{BitNet b1.58}}, where each parameter can take one of three values from the set \{-1, 0, 1\}. This variant incorporates an additional zero value to the original 1-bit BitNet, yielding a parameter representation that effectively uses 1.58 bits. \text{BitNet b1.58} preserves all the core advantages of its predecessor, including a new computation approach characterized by a near-absence of multiplication operations in matrix calculations—allowing for aggressive optimization. Energy usage remains consistent with the original BitNet, and this variant greatly surpasses FP16-based LLMs in terms of memory efficiency, throughput, and latency. Moreover, \text{BitNet b1.58} introduces two new benefits: First, it enhances modeling power due to explicit feature filtering, enabled by the inclusion of 0 among weight values—an addition that can meaningfully increase 1-bit LLM performance. Second, experimental results demonstrate that \text{BitNet b1.58} achieves parity with full precision (FP16) models in both perplexity and downstream task performance at the 3B parameter scale, under matched conditions of model size, token count, and other hyperparameters.

\section{BitNet b1.58}

\text{BitNet b1.58} extends the BitNet framework, a variant of the Transformer model where the conventional \emph{nn.Linear} layers are substituted with \emph{BitLinear} layers. This model is trained anew, employing 1.58-bit quantized weights and 8-bit quantized activations. In addition to the original BitNet structure, several alterations have been incorporated, which we detail below.

\paragraph{Quantization Function.} We employ an \emph{absmean} quantization approach to restrict weight values to the set \{-1, 0, +1\}. In this method, the weight matrix is first normalized by dividing each element by its mean absolute value, after which every entry is quantized to the nearest integer within the permissible range:

\begin{equation}
    \widetilde{W} = \mathrm{RoundClip}(\frac{W}{\gamma + \epsilon}, -1, 1),
\end{equation}

\begin{equation}
    \mathrm{RoundClip}(x, a, b) = \max(a, \min(b, \mathrm{round}(x))),
\end{equation}

\begin{equation}
    \gamma = \frac{1}{nm}\sum_{ij} |W_{ij}|.
\end{equation}

For activation quantization, the procedure is consistent with BitNet, except that activations are not rescaled into the $[0, Q_b]$ interval prior to applying non-linear transformations. Instead, activations are linearly mapped to $[-Q_b, Q_b]$ on a per-token basis, which eliminates the necessity for zero-point quantization. This modification simplifies implementation and system optimization, and according to our findings, yields almost no impact on model accuracy.

\paragraph{LLaMA-alike Components.}

With LLaMA~\cite{llama, llama2} emerging as the reference design for open-source large language models, our version of \text{BitNet b1.58} integrates architectural elements modeled after LLaMA. This includes employing RMSNorm~\cite{rmsnorm}, SwiGLU~\cite{swiglu}, rotary position embedding~\cite{rope}, and omitting all bias terms. Through this design, \text{BitNet b1.58} achieves compatibility with widely-used open-source platforms (such as Huggingface, vLLM~\cite{vllm}, and llama.cpp\footnote{\url{https://github.com/ggerganov/llama.cpp}}), facilitating effortless adoption.

\section{Results}

We conducted a comparative analysis between \text{BitNet b1.58} and our FP16 \text{LLaMA LLM} baseline across several model scales. Both models were trained for 100 billion tokens on the RedPajama dataset~\cite{redpajama}, guaranteeing consistency in the evaluation protocol. Zero-shot performance was assessed on a suite of language understanding benchmarks, namely ARC-Easy~\cite{arc}, ARC-Challenge~\cite{arc}, Hellaswag~\cite{hellaswag}, Winogrande~\cite{winoGrande}, PIQA~\cite{piqa}, OpenbookQA~\cite{openbookqa}, and BoolQ~\cite{boolq}. Additionally, validation perplexity scores were obtained using the WikiText2~\cite{wikitext2} and C4~\cite{c4} corpora.

For resource efficiency comparisons, both runtime GPU memory consumption and inference latency were systematically measured for \text{LLaMA LLM} and \text{BitNet b1.58}. This evaluation utilized the FasterTransformer framework\footnote{\url{https://github.com/NVIDIA/FasterTransformer}}, which provides optimized inference for large-scale models on GPUs. The experimental setup incorporated the 2-bit kernel from Ladder~\cite{ladder} for the \text{BitNet b1.58} model. We report the latency per generated output token as it dominates the overall inference cost.

Perplexity and computational resource findings for both \text{BitNet b1.58} and \text{LLaMA LLM} are provided in Table~\ref{tab:ppl}. The results indicate that, starting at the 3B parameter level, \text{BitNet b1.58} achieves comparable perplexity to the FP16 \text{LLaMA LLM}, while realizing a 2.71$\times$ increase in inference speed and a 3.55$\times$ reduction in GPU memory usage. Notably, the 3.9B \text{BitNet b1.58} instance attains a 2.4$\times$ speedup and 3.32$\times$ lower memory overhead compared to \text{LLaMA LLM} 3B, alongside a marked improvement in performance.

Table~\ref{tab:zero-shot} presents comprehensive zero-shot accuracy metrics for downstream tasks. Evaluation procedures followed the \emph{lm-evaluation-harness} pipeline\footnote{\url{https://github.com/EleutherAI/lm-evaluation-harness}}. These results reveal that as model size increases, the disparity in task performance between \text{BitNet b1.58} and \text{LLaMA LLM} diminishes. Importantly, \text{BitNet b1.58} achieves parity with full-precision baselines from the 3B model size upward. Consistent with perplexity observations, the 3.9B \text{BitNet b1.58} not only surpasses \text{LLaMA LLM} 3B in task accuracy, but also does so with lower memory and latency costs. Collectively, this establishes \text{BitNet b1.58} as a Pareto-dominant improvement over prevailing LLM benchmarks.

\paragraph{Memory and Latency}

The model was subsequently expanded to larger sizes, specifically 7B, 13B, and 70B parameters, and the associated computational costs were assessed. As depicted in Figure~\ref{fig:latency-memory-energy}, both latency and memory requirements were found to exhibit scalable trends, with performance acceleration improving as model scale increases. Notably, \text{BitNet b1.58} at the 70B scale achieves a 4.1-fold speed advantage over the \text{LLaMA LLM} reference model. This gain is attributed to the increased time demands of the \emph{nn.Linear} operation with larger models. Memory utilization demonstrates a comparable pattern, since the embedding layer is preserved in full precision and constitutes a diminishing share of the total memory footprint as the model grows. Both latency and memory were measured with a 2-bit kernel, suggesting potential for further cost reductions via future optimizations.

\paragraph{Energy}

To gauge energy efficiency, we estimated the arithmetic operations energy use for both \text{BitNet b1.58} and \text{LLaMA LLM}. The primary focus was on matrix multiplication, which dominates the computational expense in large language models. Figure~\ref{fig:energy} presents the energy consumption breakdown: the majority of computation in \text{BitNet b1.58} is executed as INT8 additions, in contrast to \text{LLaMA LLM}, which primarily involves FP16 additions and multiplications. Drawing on the energy consumption models from~\cite{energycost, pokebnn}, \text{BitNet b1.58} reduces matrix multiplication energy use by a factor of 71.4 on 7nm hardware. Further, we assessed total energy expenditure for inference on 512-token sequences. The findings indicate that, with increasing model size, \text{BitNet b1.58} displays markedly better energy efficiency relative to the FP16 \text{LLaMA LLM} baseline. This improvement is primarily because the contribution of \emph{nn.Linear} to the total computation increases in larger models, while the proportion of overhead from other layers diminishes.

\paragraph{Throughput}

To assess throughput, we evaluated both \text{BitNet b1.58} and the \text{LLaMA LLM} model, each with 70B parameters, across two 80GB A100 GPUs. Pipeline parallelism~\cite{gpipe} was utilized to ensure the \text{LLaMA LLM} 70B could fit within the memory constraints of the devices. Batch size was progressively increased until the GPUs reached their memory capacity, fixing the sequence length at 512. As summarized in Table~\ref{tab:throughput}, the 70B \text{BitNet b1.58} supports a batch size up to 11 times greater than that of the \text{LLaMA LLM}, resulting in an 8.9-fold improvement in throughput.

\textbf{\text{BitNet b1.58} establishes a novel scaling paradigm relating model efficiency to inference cost and model size.} Based on data presented in Figure~\ref{fig:latency-memory-energy} and \ref{fig:energy}, we note the following correspondences between 1.58-bit and 16-bit models:

\begin{itemize}
    \item BitNet b1.58 with 13B parameters surpasses the 3B FP16 LLM in latency, memory efficiency, and energy consumption.
    \item BitNet b1.58 with 30B parameters outperforms the 7B FP16 LLM with respect to latency, memory utilization, and energy usage.
    \item BitNet b1.58 with 70B parameters exceeds the 13B FP16 LLM for the same efficiency metrics.
\end{itemize}

\paragraph{Training with 2T Tokens}

Token count is a fundamental variable influencing LLM development. To explore the scalability of \text{BitNet b1.58} in this aspect, we trained a \text{BitNet b1.58} model using 2T tokens according to the data regimen of StableLM-3B~\cite{StableLM-3B-4E1T}, a leading open-source 3B model. Both models underwent evaluation on a suite of benchmarks including Winogrande~\cite{winoGrande}, PIQA~\cite{piqa}, SciQ~\cite{sciq}, LAMBADA~\cite{lambada}, and ARC-easy~\cite{arc}. Zero-shot accuracy, as presented in Table~\ref{tab:2t}, is reported; where tasks involved both accuracy and normalized accuracy, their average is used. For StableLM 3b at 2T tokens, results are sourced from its official technical report. The outcomes indicate that \text{BitNet b1.58} consistently delivers better performance on all evaluated tasks, providing evidence of robust generalization in 1.58-bit LLMs.

\section{Discussion and Future Work}

\textbf{1-bit Mixture-of-Experts (MoE) LLMs}

The Mixture-of-Experts (MoE) framework has emerged as a resource-efficient solution for scaling LLMs. Although it substantially decreases computational FLOPs, its practical use remains hindered by high memory requirements and substantial inter-device communication costs. Adopting 1.58-bit LLMs could effectively mitigate these obstacles. The reduced memory usage makes it feasible to deploy MoE models with fewer hardware resources. Additionally, it lessens the overhead associated with network-based activation transfers. In ideal circumstances, if a model can reside entirely within a single chip’s memory, communication overhead can be nearly eliminated.

\textbf{Native Support of Long Sequence in LLMs}

Handling extended input sequences has become increasingly essential for contemporary LLMs. One persistent difficulty during long-sequence inference is the substantial memory load introduced by KV caching. \text{BitNet b1.58} advances progress in this area by compressing activation precision from 16 bits to 8 bits, thus permitting the context window to expand twofold within equivalent hardware constraints. Furthermore, it opens possibilities for lossless compression to as little as 4 bits, or potentially even further reduction in future 1.58-bit LLM models, which is an area we intend to explore further.

\textbf{LLMs on Edge and Mobile}

Deploying 1.58-bit LLMs offers significant benefits for edge and mobile environments, which are commonly constrained by limited memory and processing capabilities. The drastic decrease in memory and power demands makes these models viable for real-time applications on such devices, unlocking functionalities that previously were unfeasible. This not only broadens the spectrum of feasible use cases for edge and mobile computing but also fosters innovation in LLM-driven applications. Notably, 1.58-bit LLMs are particularly well-suited to CPUs, which dominate edge and mobile processing; this enables \text{BitNet b1.58} to execute efficiently and capitalize on the computational resources typically available in these platforms.

\textbf{New Hardware for 1-bit LLMs}

Recent advancements, exemplified by~Groq\footnote{\url{https://groq.com/}}, highlight the potential of designing specialized hardware architectures (e.g., LPUs) targeted at large language models. Building upon these developments, we propose dedicated efforts toward engineering new hardware and systems customized for the computational model introduced by 1-bit LLMs, leveraging the distinctive advantages of BitNet~\cite{bitnet}’s paradigm.

\end{document}